\documentclass[10pt,journal,compsoc]{IEEEtran}

\usepackage{amsmath}
\usepackage{amssymb}
\usepackage{algorithm}
\usepackage{algpseudocode}
\usepackage{booktabs}
\usepackage{array}
\usepackage{multirow}
\usepackage{graphicx}
\usepackage[hidelinks]{hyperref}
\usepackage{url}
\usepackage{cite}
\usepackage{microtype}
\usepackage{placeins}

% IEEEtran column width is ~252pt; set tabcolsep accordingly
\setlength{\tabcolsep}{4pt}

\algnewcommand\algorithmicinput{\textbf{Input:}}
\algnewcommand\algorithmicoutput{\textbf{Output:}}
\algnewcommand\Input{\item[\algorithmicinput]}
\algnewcommand\Output{\item[\algorithmicoutput]}

\begin{document}

\title{Zero-Shot Multi-Core Process Scheduling via Offline
  Monte Carlo Tree Search with Contrastive Value Learning}

\author{Sanusi~Mustapha~Babansoro~and~Abrha~Gebriye~Embafresu%
\IEEEcompsocitemizethanks{
\IEEEcompsocthanksitem S.\,M.~Babansoro and A.\,G.~Embafresu are with
the School of Computer Science and Engineering, University of Electronic
Science and Technology of China (UESTC), Chengdu 611731, China.
E-mail: 202524080114@std.uestc.edu.cn (corresponding author).
}}

\markboth{IEEE Transactions on Parallel and Distributed Systems}{}

\IEEEtitleabstractindextext{%
\begin{abstract}
Process scheduling on symmetric multi-processing (SMP) systems is a
fundamental resource management challenge whose difficulty grows with core
count: work-stealing and load-balancing heuristics improve throughput but
cannot account for process priority or prevent starvation, while existing
reinforcement learning (RL) schedulers require live environment interaction
during training and must be retrained whenever the core count changes.
This paper proposes Retrospective MCTS (R-MCTS), an offline scheduling
agent that addresses both limitations. R-MCTS improves its scheduling
policy by running UCT tree search over stored failed trajectories without
any live system interaction, then deploys the learned policy across SMP
systems with varying core counts without retraining. A priority-aware
value estimator ($6{\to}24{\to}12{\to}1$ MLP) trained with contrastive
offline learning guides the search; an approximate nearest-neighbour
policy table enables runtime generalisation to previously unseen scheduler
states. We evaluate R-MCTS on three workload types (CPU-bound, I/O-bound,
mixed), a Google cluster trace derived workload, and four SMP
configurations ($N\!=\!1,2,4,8$ cores), reporting 10 experiments with 15
to 20 independent trials each. On single-core mixed workloads, R-MCTS
achieves turnaround within 1.9\% of the SJF-True theoretical ceiling
($p{=}0.667$), outperforms SJF-Est by 23.8\% ($p{<}0.001$), and reduces
short-process wait time by 90.0\% under starvation stress. On the Google
trace workload R-MCTS achieves the lowest starvation count and highest
CPU utilisation (90.8\%). In the multi-core transfer study, R-MCTS trained
exclusively on single-core trajectories achieves statistically
indistinguishable performance from FCFS on $N\!=\!2,4,8$ cores
($p{>}0.70$), demonstrating zero-shot core-count generalisation. An
ablation study confirms the anti-starvation mechanism contributes 5.8\%
and UCT search contributes 0.8\% of the single-core turnaround gain.
\end{abstract}

\begin{IEEEkeywords}
Process scheduling, Monte Carlo tree search, offline reinforcement
learning, symmetric multiprocessing, starvation prevention,
resource management.
\end{IEEEkeywords}}

\maketitle

\IEEEdisplaynontitleabstractindextext
\IEEEpeerreviewmaketitle

% ─────────────────────────────────────────────────────────────────────────────
\section{Introduction}
\label{sec:intro}
% ─────────────────────────────────────────────────────────────────────────────

\IEEEPARstart{P}{rocess} scheduling is among the most consequential
resource management decisions in modern computing systems.
On symmetric multi-processing (SMP) platforms---which now constitute the
default architecture from smartphone SoCs to server chips---the scheduler
must simultaneously maximise throughput, maintain response-time fairness
across heterogeneous priority classes, and prevent indefinite starvation
of low-priority processes when high-priority work is abundant.
Classical algorithms make clear tradeoffs:
FCFS is starvation-free but causes convoy effects;
non-preemptive SJF minimises mean wait time but requires exact burst
knowledge unavailable at runtime;
Round-Robin provides fairness at the cost of frequent context switches;
MLFQ approximates SJF without burst knowledge by demoting CPU-bound
processes across priority queues.
None of these policies adapt as workloads evolve, and none were designed
for multi-core coordination: ported naively to $N$ cores they share a
global ready queue but apply the same per-process selection logic,
ignoring the load-imbalance signal that emerges when cores compete for
a finite set of ready processes~\cite{stallings,silberschatz}.

\subsection{Problem Statement}

Two limitations constrain the applicability of existing learned schedulers
to production SMP systems.
First, online RL approaches~\cite{sundqn,hanlearning,linnos} require live
environment interaction during training, injecting exploration decisions
into the actual scheduler and creating safety and stability risks that
are unacceptable in deployed systems.
Second, schedulers trained on a fixed core count must be retrained from
scratch when the hardware configuration changes---a significant practical
burden for operators managing heterogeneous fleets.

This paper addresses both limitations through a single unified mechanism:
\emph{offline retrospective improvement}. The central question is:
\emph{can a scheduling agent trained exclusively on single-core
trajectories generalise, without retraining, to SMP systems with two,
four, or eight cores?}
We answer affirmatively through R-MCTS, whose per-process feature
representation is core-count independent by design.

\subsection{Contributions}

\begin{enumerate}
  \item \textbf{Offline MCTS for SMP scheduling.}
    R-MCTS is the first scheduler to apply UCT-guided offline tree search
    to the multi-core process scheduling problem.
    Training requires only stored single-core trajectories; no live
    environment access is needed at any point.

  \item \textbf{Zero-shot core-count generalisation.}
    We prove empirically that R-MCTS trained on $N{=}1$ core achieves
    statistically indistinguishable turnaround from baseline schedulers
    on $N{=}2,4,8$ cores without retraining, a property we term
    \emph{core-count transfer}.

  \item \textbf{Comprehensive evaluation.}
    Ten experiments spanning synthetic workloads, Google cluster
    traces~\cite{reiss2011google}, and four SMP configurations,
    including an ablation study and full Welch $t$-test significance
    analysis across 15--20 independent trials per condition.

  \item \textbf{Open simulator.}
    A fully instrumented discrete-event OS simulator implementing the
    five-state process model, SMP scheduling, and I/O burst cycles,
    validated against theoretical predictions.
\end{enumerate}

% ─────────────────────────────────────────────────────────────────────────────
\section{Related Work}
\label{sec:related}
% ─────────────────────────────────────────────────────────────────────────────

\subsection{Classical OS Scheduling}

Stallings~\cite{stallings} and Silberschatz et~al.~\cite{silberschatz}
provide canonical treatments of FCFS, SJF, Round-Robin, and MLFQ.
The Linux Completely Fair Scheduler (CFS)~\cite{silberschatz} used a
red-black tree tracking virtual runtime, providing nanosecond-granularity
fairness but no explicit starvation bound.
Linux 6.6 replaced CFS with EEVDF (Earliest Eligible Virtual Deadline
First), which schedules based on virtual deadlines rather than virtual
runtime, improving worst-case latency~\cite{silberschatz}.
Both CFS and EEVDF are engineered heuristics; neither adapts its
parameters from historical scheduling data.

\subsection{Reinforcement Learning for Scheduling}

Deep RL has been applied to cluster and resource scheduling with notable
success. Mao et~al.~\cite{mao2016} pioneered RL-based cluster resource
management; Mao et~al.~\cite{decima} extended this to data-processing
clusters with graph neural networks.
For OS-level scheduling, Sun et~al.~\cite{sundqn} applied Double DQN to
task scheduling and outperformed static policies on dynamic workloads.
Han et~al.~\cite{hanlearning} applied policy gradient methods to Linux
CPU scheduler parameter tuning on real hardware.
Hao et~al.~\cite{linnos} embedded a light neural network in the Linux
storage I/O path, validating sub-millisecond learned decisions within the
OS kernel.
All of these approaches require online interaction during training.
For multi-core and distributed scheduling, Rjoub et~al.~\cite{rjoubccpe}
applied DQN and DRL-LSTM to cloud task scheduling, and
Barrett et~al.~\cite{barrettccpe} automated cloud resource allocation
through RL.
Pang et~al.~\cite{pangccpe} applied sparse-reward DRL to cloud
manufacturing scheduling.
A critical review by Jalali Khalil Abadi et~al.~\cite{jalali2024} covers
the broader landscape of RL-based scheduling in distributed systems.
None of these works study core-count transfer or offline training.

\subsection{Offline Reinforcement Learning}

Offline RL~\cite{levine2020} trains from a fixed dataset without
environment access, addressing distribution shift through trajectory
quality weighting.
Q-Transformer~\cite{qtransformer} extends offline RL to autoregressive
architectures.
Direct Preference Optimisation~\cite{dpo} and AlphaLLM~\cite{alphallm}
use contrastive preference pairs for offline value learning.
R-MCTS adapts the preference-pair framework to process-level scheduling:
improved trajectories from UCT search are paired with failed ones,
and the value estimator is trained to rank candidate processes rather
than score whole episodes.

\subsection{Monte Carlo Tree Search}

UCT was formalised by Kocsis and Szepesv\'{a}ri~\cite{uct} and surveyed
by Browne et~al.~\cite{mcts_survey}.
AlphaGo~\cite{alphago} and MuZero~\cite{muzero} demonstrated that UCT
combined with learned value functions produces superhuman policies.
MCTS-EP~\cite{mctsep} applied offline MCTS to embodied task planning.
R-MCTS is the first application of offline UCT to OS process scheduling
and the first to study its behaviour on SMP systems.

% ─────────────────────────────────────────────────────────────────────────────
\section{System Design}
\label{sec:design}
% ─────────────────────────────────────────────────────────────────────────────

\subsection{Single-Core Simulator}

The simulator models the five-state process lifecycle~\cite{stallings}:
\textsc{New}, \textsc{Ready}, \textsc{Running}, \textsc{Blocked},
\textsc{Exit}.
Each process is represented by a PCB (Table~\ref{tab:pcb}).
Processes alternate between CPU execution and I/O blocking phases.
Context switch overhead is one tick per switch.
Figure~\ref{fig:architecture} shows the complete system architecture.

\begin{table}[!t]
\centering
\caption{Process Control Block (PCB) Field Specification}
\label{tab:pcb}
\renewcommand{\arraystretch}{1.1}
\begin{tabular}{@{}l l p{4.2cm}@{}}
\toprule
\textbf{Field} & \textbf{Type} & \textbf{Role} \\
\midrule
\texttt{pid}            & int  & Unique process identifier \\
\texttt{burst\_time}    & int  & Total CPU execution time required \\
\texttt{remaining}      & int  & Remaining burst ticks \\
\texttt{priority}       & int  & Priority (1\,=\,highest, 5\,=\,lowest) \\
\texttt{arrival\_time}  & int  & Tick at which process enters \textsc{New} \\
\texttt{state}          & enum & Five-state process state \\
\texttt{io\_burst}      & int  & Duration of each I/O blocking phase \\
\texttt{wait\_time}     & int  & Ticks accumulated in Ready queue \\
\texttt{response\_time} & int  & Ticks from arrival to first dispatch \\
\texttt{finish\_time}   & int  & Tick at which process reaches Exit \\
\bottomrule
\end{tabular}
\end{table}

\begin{figure*}[!t]
  \centering
  \includegraphics[width=\textwidth]{OS ARCHITECTURE-1.png}
  \caption{R-MCTS system architecture. Stage~A: OS process simulator
    implementing the five-state model. Stage~B: trajectory buffer storing
    all scheduling decisions and outcomes. Stage~C: value estimator $V_\theta$
    trained with contrastive offline preference learning.
    Stage~D: UCT tree search operating on stored failed trajectories.
    Stage~E: policy improvement loop feeding updated policy back to Stage~A.
    Baselines (FCFS, SJF-True, SJF-Est, Round-Robin, MLFQ) share Stage~A.}
  \label{fig:architecture}
\end{figure*}

\subsection{SMP Extension}

The SMP simulator extends the single-core design with $N$ independent
CPU cores sharing one ready queue and one blocked queue
(Algorithm~\ref{alg:smp}).
This follows Linux's global runqueue model~\cite{silberschatz}, enabling
clear analysis of the scheduling decision in isolation from
work-stealing mechanics.
At each tick, every idle core calls the scheduler function independently;
the scheduler sees the current ready queue state and selects one process.
Per-core utilisation and load-imbalance ratio (standard deviation of
per-core utilisation across all cores) are recorded alongside the
standard single-core metrics.

\begin{algorithm}[!t]
\caption{SMP Simulator: one tick}
\label{alg:smp}
\begin{algorithmic}[1]
\Input{$N$ cores $\mathcal{C}$, shared $\mathcal{RQ}$,
  $\mathcal{BQ}$, scheduler $\pi$}
\State Admit newly arrived processes to $\mathcal{RQ}$
\State Unblock I/O-complete processes: $\mathcal{BQ}.\text{tick}()$
\For{each idle core $c \in \mathcal{C}$}
  \If{$\mathcal{RQ} \ne \emptyset$}
    \State $p^* \leftarrow \pi(\mathcal{RQ}, c)$;\;
           $\mathcal{RQ}.\text{remove}(p^*)$;\;
           $c.\text{assign}(p^*)$
  \EndIf
\EndFor
\For{each core $c \in \mathcal{C}$}
  \State $c.\text{tick}()$\;
         (handles context switches, I/O preemption, process completion)
\EndFor
\State Increment $\text{wait\_time}$ for all processes in $\mathcal{RQ}$
\end{algorithmic}
\end{algorithm}

\subsection{Mathematical Notation}

Table~\ref{tab:notation} consolidates all symbols used throughout.

\begin{table}[!t]
\centering
\caption{Mathematical Notation}
\label{tab:notation}
\renewcommand{\arraystretch}{1.1}
\begin{tabular}{@{}l p{5.5cm}@{}}
\toprule
\textbf{Symbol} & \textbf{Meaning} \\
\midrule
$\tau$, $\tau^+$, $\tau^-$ & Trajectory, improved, original (failed) \\
$R_{\text{pw}}(\tau)$       & Priority-weighted turnaround of $\tau$ \\
$T_i$, $w(p_i)$             & Turnaround and priority weight of process $i$ \\
$\mu_B$, $\sigma_B$         & Mean and std.\ dev.\ of $R_{\text{pw}}$ in buffer \\
$V_\theta$                  & Value estimator MLP with weights $\theta$ \\
$f_i$                       & 6-dim feature vector for process $i$ \\
$\mathbf{s}$                & 8-dim queue state feature vector \\
$Q(s,a)$, $N(s)$, $N(s,a)$ & UCT value, state and action visit counts \\
$c$                         & UCT exploration constant ($\sqrt{2}$) \\
$\Pi$                       & Approximate nearest-neighbour policy table \\
$\delta$                    & Anti-starvation threshold (ticks) \\
$N_c$                       & Number of CPU cores \\
$\sigma_U$                  & Load imbalance (std.\ dev.\ of per-core utilisation) \\
\bottomrule
\end{tabular}
\end{table}

% ─────────────────────────────────────────────────────────────────────────────
\section{R-MCTS Agent}
\label{sec:agent}
% ─────────────────────────────────────────────────────────────────────────────

\subsection{Value Estimator}

$V_\theta$ is a three-layer MLP ($6{\to}24{\to}12{\to}1$) with ReLU
hidden activations and a sigmoid output, mapping a 6-dimensional process
feature vector $f_i$ to a quality score in $[0,1]$.
Input features are: normalised remaining burst, priority encoding,
normalised wait time, remaining-to-burst ratio, starvation signal, and
normalised queue length.
Training uses contrastive process-level labels: for each UCT-improved
decision the chosen process receives label $y{=}1.0$; all other processes
in the same queue receive $y{=}0.2$.
This teaches $V_\theta$ to \emph{rank} candidates, not score episodes.
The priority-weighted return and normalised training target are:
\begin{align}
  R_{\text{pw}}(\tau) &= \frac{\sum_i w(p_i)\,T_i}{\sum_i w(p_i)}
  \label{eq:rpw}\\[4pt]
  y &= \operatorname{clip}\!\Bigl(
    1 - \tfrac{R_{\text{pw}} - \mu_B}{3\sigma_B + \varepsilon},\;0,\;1
  \Bigr)
  \label{eq:target}
\end{align}
with $w\!=\!5.0$ for priority~1 decreasing to $w\!=\!1.0$ for
priority~5, and $\varepsilon\!=\!10^{-8}$.
Training minimises MSE with $\alpha\!=\!0.01$ over 40 epochs per
iteration (Algorithm~\ref{alg:training}).

\begin{algorithm}[!t]
\caption{Value Estimator Training (Contrastive)}
\label{alg:training}
\begin{algorithmic}[1]
\Input{Dataset $\mathcal{D}{=}\{(f_i,y_i)\}$,
  $\alpha{=}0.01$, $E{=}40$, batch $B{=}32$}
\Output{Updated weights $\theta$}
\State Initialise $\theta$ with Xavier uniform initialisation
\For{each epoch $e = 1,\ldots,E$}
  \For{each mini-batch $\{(f_i, y_i)\}_{i=1}^{B}$}
    \State $\hat{y}_i \leftarrow V_\theta(f_i)$;\;
           $\mathcal{L} = \frac{1}{B}\sum_i(\hat{y}_i-y_i)^2$;\;
           $\theta \leftarrow \theta - \alpha\nabla_\theta\mathcal{L}$
  \EndFor
\EndFor
\end{algorithmic}
\end{algorithm}

\subsection{UCT Tree Search and Offline Improvement}

For each failed trajectory in the replay buffer R-MCTS identifies the
failure point (decision with highest average ready-queue wait time), then
runs $K{=}60$ UCT simulations from that state using $V_\theta$:
\begin{equation}
  \text{UCT}(s,a) = Q(s,a) + c\sqrt{\frac{\ln N(s)}{N(s,a)}}
  \label{eq:uct}
\end{equation}
Improved trajectories are paired with failed ones; the preference weight
\begin{equation}
  w(\tau) = \frac{\max(0,\;R_{\text{pw}}(\tau^+)-R_{\text{pw}}(\tau^-))}
    {R_{\text{pw}}(\tau^-)+\varepsilon}
  \label{eq:weight}
\end{equation}
filters out non-improvements.
Algorithms~\ref{alg:improvement}--\ref{alg:policy} give the full
offline loop and runtime scheduling procedure.

\begin{algorithm}[!t]
\caption{R-MCTS Offline Improvement Loop}
\label{alg:improvement}
\begin{algorithmic}[1]
\Input{Buffer $\mathcal{B}$, $V_\theta$, $K{=}60$, $M{=}30$}
\Output{Policy table $\Pi$, updated $V_\theta$}
\State Collect $M$ warmup episodes; store in $\mathcal{B}$
\For{iteration $t = 1,\ldots,10$}
  \State Add 10 new episodes to $\mathcal{B}$
  \For{each failed $\tau^- \in \mathcal{B}$}
    \State $s^* \leftarrow$ highest-wait decision in $\tau^-$
    \State Run $K$-sim UCT from $s^*$; obtain $\tau^+$, $p^*$
    \If{$w(\tau) > 0$}
      \State Insert $(s^*,\,f(p^*))$ into $\Pi$
    \EndIf
  \EndFor
  \State Retrain $V_\theta$ on preference dataset
\EndFor
\end{algorithmic}
\end{algorithm}

\begin{algorithm}[!t]
\caption{R-MCTS Runtime Scheduling}
\label{alg:policy}
\begin{algorithmic}[1]
\Input{Ready queue $\mathcal{RQ}$, clock $t$, $\Pi$, $\delta{=}25$}
\Output{Selected process $p^*$}
\State $\hat{f} \leftarrow \Pi.\text{approxLookup}(\mathbf{s})$
\If{$\hat{f} \ne$ \textbf{None}}
  \State \Return $\arg\min_{p \in \mathcal{RQ}} \|f(p)-\hat{f}\|$
\EndIf
\State Apply emergency priority boost to all $p$ with $\text{wait}(p)\ge\delta$
\State \Return $\arg\max_{p \in \mathcal{RQ}} V_\theta(f(p))$
\end{algorithmic}
\end{algorithm}

\subsection{Core-Count Transfer Property}

A key design property of Algorithm~\ref{alg:policy} is that the
scheduler function takes only per-process features ($f_i$) and queue
state ($\mathbf{s}$) as inputs; neither the number of cores $N_c$ nor
any inter-core state appears in the feature vector.
This means the same policy table $\Pi$ and value estimator $V_\theta$
trained on $N_c{=}1$ can be deployed on $N_c{=}2,4,8$ without
modification.
At each scheduling point, each idle core independently queries the same
policy, observing the current shared ready queue.
Section~\ref{sec:exp9} validates this property experimentally.

\subsection{Hyperparameter Configuration}

All hyperparameters are listed in Table~\ref{tab:hyperparams} and held
fixed across all experiments.

\begin{table}[!t]
\centering
\caption{R-MCTS Hyperparameter Configuration}
\label{tab:hyperparams}
\renewcommand{\arraystretch}{1.1}
\begin{tabular}{@{}l l p{3.8cm}@{}}
\toprule
\textbf{Symbol} & \textbf{Value} & \textbf{Description} \\
\midrule
$K$           & 60         & UCT simulations per failure point \\
$M$           & 30         & Warmup episodes \\
$\delta$      & 25         & Starvation wait threshold (ticks) \\
$\alpha$      & 0.01       & Value estimator learning rate \\
$E$           & 40         & Training epochs per iteration \\
$c$           & $\sqrt{2}$ & UCT exploration constant \\
$w_1$--$w_5$ & 5.0--1.0   & Priority weights for levels 1--5 \\
ANN threshold & 0.88       & Cosine similarity threshold \\
\bottomrule
\end{tabular}
\end{table}

% ─────────────────────────────────────────────────────────────────────────────
\section{Experimental Design}
\label{sec:experiments}
% ─────────────────────────────────────────────────────────────────────────────

\subsection{Development Environment}

All experiments run on Windows~11 with Python~3.11.2, NumPy~1.21+,
Matplotlib~3.4+, SciPy, and our discrete-event simulator
(no kernel modification, no GPU required).
Full source code is available at~\cite{rmcts_github}.
Total runtime: approximately 280 minutes for all ten experiments.

\subsection{Workload Types}

Three synthetic workload types follow standard OS workload
classification~\cite{stallings} (Table~\ref{tab:workloads}).
A fourth workload derives statistical properties from the Google Cluster
Usage Trace~2011~\cite{reiss2011google}: burst times follow a
lognormal distribution (median $\approx$5 ticks, mean $\approx$35 ticks)
matching the heavy-tailed task duration distribution reported by
Minet et~al.~\cite{minet2018}; priorities follow Google's
0--11 scale mapped to our 1--5 scale; arrivals follow a Poisson process.

\begin{table}[!t]
\centering
\caption{Synthetic Workload Parameters (20 Processes/Trial)}
\label{tab:workloads}
\renewcommand{\arraystretch}{1.1}
\begin{tabular}{@{}l c c c@{}}
\toprule
\textbf{Parameter} & \textbf{CPU-Bound} & \textbf{I/O-Bound} & \textbf{Mixed} \\
\midrule
CPU phase (ticks) & 8--20 & 2--5  & 2--18 \\
I/O burst (ticks) & 1--3  & 5--14 & 1--12 \\
I/O phases        & 1--2  & 2--5  & 1--4 \\
Trials            & 20    & 20    & 20 \\
\bottomrule
\end{tabular}
\end{table}

\subsection{Baseline Schedulers}

Six schedulers are compared in single-core experiments
(Table~\ref{tab:schedulers}); MLFQ is excluded from multi-core
experiments since its per-queue demotion logic does not extend cleanly
to the shared-queue SMP model used here.

\begin{table}[!t]
\centering
\caption{Baseline Scheduler Descriptions}
\label{tab:schedulers}
\renewcommand{\arraystretch}{1.1}
\begin{tabular}{@{}l l p{4.0cm}@{}}
\toprule
\textbf{Scheduler} & \textbf{Category} & \textbf{Description} \\
\midrule
FCFS      & Baseline  & Arrival-order; no priority. \\
SJF-True  & Ceiling   & Exact burst times; unrealisable. \\
SJF-Est   & Realistic & Exponential averaging ($\alpha{=}0.5$). \\
RR        & Baseline  & Preemptive, quantum\,$=3$ ticks. \\
MLFQ      & Baseline  & 3 levels, quanta 3/6/12, aging 20 ticks. \\
R-MCTS    & Proposed  & See Section~\ref{sec:agent}. \\
\bottomrule
\end{tabular}
\end{table}

\subsection{Performance Metrics}

\begin{itemize}
  \item \textbf{Avg.\ turnaround time}: mean of
    $(\text{finish\_time}-\text{arrival\_time})$.
  \item \textbf{Avg.\ wait time}: mean ticks spent in ready queue.
  \item \textbf{Starvation count}: processes with wait time exceeding
    $3{\times}$ burst time.
  \item \textbf{CPU utilisation}: $(T_{\text{total}}-T_{\text{idle}})/T_{\text{total}}$,
    averaged across cores for SMP.
  \item \textbf{Load imbalance} ($\sigma_U$): standard deviation of
    per-core utilisation; smaller is better.
\end{itemize}

\subsection{Statistical Significance}

All two-scheduler comparisons use the Welch $t$-test (unequal variance);
$p{<}0.05$ is the significance threshold.
Effect sizes are reported as percentage improvement in mean turnaround.

% ─────────────────────────────────────────────────────────────────────────────
\section{Single-Core Results (Experiments 1--7)}
\label{sec:single}
% ─────────────────────────────────────────────────────────────────────────────

\subsection{Value Estimator Learning Curve (Exp.\ 1)}

Training loss decreases from 0.0544 to 0.0392 (28.0\% reduction) over
ten iterations on synthetic mixed workloads, and from 0.0476 to 0.0445
(6.6\% reduction) on the Google trace workload, reflecting the greater
difficulty of the heavy-tailed burst distribution.
The policy table grows from 0 to 150 entries on synthetic workloads
(Figure~\ref{fig:learning}).

\begin{figure}[!t]
  \centering
  \includegraphics[width=\linewidth]{fig1_learning_1col.png}
  \caption{R-MCTS learning curve (synthetic mixed workload, 10 iterations):
    (a) MSE training loss decreases 28.0\%;
    (b) policy table grows 0$\to$150 entries.}
  \label{fig:learning}
\end{figure}

\subsection{Main Benchmark: Mixed Workload (Exp.\ 2)}

Table~\ref{tab:main} reports single-core performance on the mixed
workload (20 trials, 20 processes each).
R-MCTS achieves $88.71{\pm}11.82$ ticks---within 1.9\% of the SJF-True
ceiling, a gap that is not statistically significant ($p{=}0.667$).
R-MCTS outperforms SJF-Est by 23.8\% ($p{<}0.001$),
Round-Robin by 42.3\% ($p{<}0.001$), and MLFQ by 34.7\% ($p{<}0.001$).
R-MCTS and FCFS achieve identical CPU utilisation (97\%), confirming
that improved turnaround carries no utilisation cost.
Figure~\ref{fig:mixed} visualises these comparisons.

\begin{table*}[!t]
\centering
\caption{Single-Core Benchmark: Mixed Workload (20 Trials, 20 Processes).
  Mean $\pm$ SD in ticks.}
\label{tab:main}
\renewcommand{\arraystretch}{1.15}
\begin{tabular}{@{}l r r r r@{}}
\toprule
\textbf{Scheduler} & \textbf{Turnaround} & \textbf{Wait}
  & \textbf{Starv.} & \textbf{CPU\%} \\
\midrule
FCFS               & $116.36\pm14.22$ & $100.22\pm13.46$ & 16.6 & 97 \\
SJF-True$^\dagger$ & $ 87.03\pm12.04$ & $ 70.89\pm11.17$ & 14.6 & 97 \\
SJF-Est            & $116.36\pm14.22$ & $100.22\pm13.46$ & 16.6 & 97 \\
Round-Robin        & $153.68\pm24.08$ & $137.55\pm23.21$ & 19.6 & 95 \\
MLFQ               & $135.80\pm18.33$ & $119.66\pm17.52$ & 19.6 & 96 \\
\textbf{R-MCTS}$^*$& $\mathbf{88.71\pm11.82}$ & $\mathbf{72.58\pm10.94}$
  & \textbf{15.7} & \textbf{97} \\
\bottomrule
\multicolumn{5}{@{}l@{}}{\footnotesize
  $\dagger$ Theoretical ceiling (exact burst times, unrealisable).}\\
\multicolumn{5}{@{}l@{}}{\footnotesize
  $^*$ Proposed. Welch $t$-test $p{<}0.001$ vs all practical baselines.}
\end{tabular}
\end{table*}

\begin{figure*}[!t]
  \centering
  \includegraphics[width=\linewidth]{fig2_comparison_2col.png}
  \caption{Single-core mixed workload comparison: (a)~turnaround,
    (b)~wait time, (c)~starvation count. Dots are means; caps show
    $\pm 1$ SD. R-MCTS bar outlined in black.}
  \label{fig:mixed}
\end{figure*}

\subsection{Starvation Stress Test (Exp.\ 3)}

Under a deliberately adversarial workload (one long-burst process,
burst 40--60 ticks, followed by four short-burst processes, burst 2--5
ticks), FCFS forces short processes to wait $55.64{\pm}6.10$ ticks.
R-MCTS reduces this to $5.54{\pm}0.95$ ticks---a \textbf{90.0\%
reduction}---approaching SJF-True ($4.44{\pm}0.89$ ticks).
Figure~\ref{fig:starvation} illustrates the results.

\begin{figure*}[!t]
  \centering
  \includegraphics[width=\linewidth]{fig3_starvation_2col.png}
  \caption{Starvation stress test (20 trials):
    (a) short-process avg.\ wait time;
    (b) overall avg.\ turnaround.
    R-MCTS reduces short-process wait by 90.0\% vs FCFS.}
  \label{fig:starvation}
\end{figure*}

\subsection{Scalability (Exp.\ 4)}

Figure~\ref{fig:scalability} shows turnaround vs process count
$N{=}5,10,20,30,50$ on the single-core mixed workload.
R-MCTS stays within 1.4\% of SJF-True at $N{=}50$, confirming that
the offline policy generalises under increasing load.

\begin{figure}[!t]
  \centering
  \includegraphics[width=\linewidth]{fig4_scalability_1col.png}
  \caption{Single-core scalability: avg.\ turnaround vs $N$
    (mixed workload, 10 trials per $N$).
    R-MCTS tracks SJF-True closely at all scales.}
  \label{fig:scalability}
\end{figure}

\subsection{Ablation Study (Exp.\ 5)}

Table~\ref{tab:ablation} decomposes R-MCTS into three variants.
The anti-starvation mechanism contributes 5.8\% turnaround improvement;
UCT offline search contributes an additional 0.8\%.
The heuristic fallback (priority-weighted SJF) carries the majority of
scheduling quality, confirming that the offline search adds targeted
refinement at historically difficult states.

\begin{table}[!t]
\centering
\caption{Ablation Study (Mixed Workload, 20 Processes, 15 Trials)}
\label{tab:ablation}
\renewcommand{\arraystretch}{1.1}
\footnotesize
\begin{tabular}{@{}l r r r@{}}
\toprule
\textbf{Variant} & \textbf{Turn.} & \textbf{Wait} & \textbf{$\Delta$Full} \\
\midrule
R-MCTS-Full      & $83.35\pm10.40$ & $67.61\pm9.81$  & baseline \\
R-MCTS-Heuristic & $83.98\pm10.81$ & $68.23\pm10.22$ & $+$0.8\% \\
R-MCTS-NoStarv   & $88.21\pm10.79$ & $72.47\pm10.20$ & $+$5.8\% \\
FCFS (ref.)      & $104.30\pm14.28$& $88.56\pm13.69$ & $+$25.1\% \\
\bottomrule
\end{tabular}
\vspace{2pt}
\footnotesize
Heuristic = priority-weighted SJF + anti-starvation, no UCT.\\
NoStarv = UCT + ANN matching, starvation threshold removed.
\end{table}

\begin{figure*}[!t]
  \centering
  \includegraphics[width=\linewidth]{fig5_ablation_2col.png}
  \caption{Ablation study: turnaround and wait for each R-MCTS variant.
    Dotted line marks FCFS reference level.}
  \label{fig:ablation}
\end{figure*}

\subsection{Cross-Workload Results (Exp.\ 6)}

Figure~\ref{fig:breakdown} shows R-MCTS achieves the lowest starvation
count in all three synthetic workload types.
Figure~\ref{fig:responsetime} shows Round-Robin achieves the lowest
response time on I/O-bound workloads (22.69 ticks) due to preemption;
R-MCTS accepts higher response time (26.10 ticks) in exchange for
substantially lower turnaround (22.0\% lower on I/O-bound, 42.3\% on
mixed) and the 90.0\% starvation reduction.

\begin{figure*}[!t]
  \centering
  \includegraphics[width=\linewidth]{fig6_workloads_2col.png}
  \caption{Performance across all three synthetic workload types:
    (a)~turnaround; (b)~starvation count.
    R-MCTS and SJF-True occupy the outermost positions in both panels.}
  \label{fig:breakdown}
\end{figure*}

\begin{figure*}[!t]
  \centering
  \includegraphics[width=\linewidth]{fig7_response_2col.png}
  \caption{Average response time across workload types.
    Round-Robin (gray) achieves lowest response via preemption;
    R-MCTS (blue) trades response time for near-optimal turnaround
    and starvation control.}
  \label{fig:responsetime}
\end{figure*}

\subsection{Google Cluster Trace Validation (Exp.\ 7)}
\label{sec:exp7}

Figure~\ref{fig:google} shows single-core results on the
Google-trace-derived workload (20 trials, 20 processes each).
R-MCTS outperforms Round-Robin by 35.3\% ($p{<}0.001$),
MLFQ by 33.6\% ($p{=}0.002$), and achieves the highest CPU utilisation
(90.8\%) of all schedulers.
The gap to SJF-True is larger here (50.4\%) than on synthetic workloads
because the heavy-tailed lognormal burst distribution makes burst
prediction harder---consistent with published analyses of Google
trace statistics~\cite{minet2018}.
Table~\ref{tab:google_stats} reports full significance results.

\begin{figure*}[!t]
  \centering
  \includegraphics[width=\linewidth]{fig8_google_trace_2col.png}
  \caption{Google cluster trace workload validation (20 trials,
    20 processes/trial, single core):
    (a)~turnaround; (b)~wait time; (c)~starvation count.}
  \label{fig:google}
\end{figure*}

\begin{table}[!t]
\centering
\caption{Statistical Significance: R-MCTS vs.\ Baselines
  (Welch $t$-Test, Google Trace, Single Core, 20 Trials)}
\label{tab:google_stats}
\renewcommand{\arraystretch}{1.1}
\footnotesize
\begin{tabular}{@{}l r r r l@{}}
\toprule
\textbf{Comp.} & \textbf{R-MCTS} & \textbf{Base.}
  & \textbf{Impr.} & \textbf{$p$} \\
\midrule
vs FCFS        & 545.79 & 606.41 & $+$10.0\% & 0.236\,ns \\
vs SJF-True    & 545.79 & 362.83 & $-$50.4\% & $<$0.001 \\
vs SJF-Est     & 545.79 & 562.53 & $+$3.0\%  & 0.745\,ns \\
vs Round-Robin & 545.79 & 843.86 & $+$35.3\% & $<$0.001 \\
vs MLFQ        & 545.79 & 822.35 & $+$33.6\% & 0.002 \\
\bottomrule
\end{tabular}
\vspace{2pt}
\footnotesize ns = not significant at $p{<}0.05$.
\end{table}

% ─────────────────────────────────────────────────────────────────────────────
\section{Multi-Core Results (Experiments 8--10)}
\label{sec:multicore}
% ─────────────────────────────────────────────────────────────────────────────

\subsection{SMP Scalability (Exp.\ 8)}
\label{sec:exp8}

Figure~\ref{fig:multicore} shows turnaround, starvation, and CPU
utilisation as core count increases from $N_c{=}1$ to $N_c{=}8$
on the mixed workload (15 trials, 20 processes each).
All schedulers benefit from additional cores: turnaround halves from
$N_c{=}1$ to $N_c{=}2$ and halves again to $N_c{=}4$, consistent with
Amdahl's law for a predominantly parallel workload.
R-MCTS achieves the \textbf{lowest load imbalance at $N_c{=}2$}
($\sigma_U{=}0.024$ vs FCFS 0.031, Round-Robin 0.039),
meaning its priority-aware value estimator distributes work more evenly
across cores than FCFS's arrival-order selection.
At $N_c{=}8$, with only 20 processes, the system is lightly loaded and
all schedulers converge; differences become negligible, which is
the correct and expected behaviour.
Starvation counts decrease monotonically with core count for all
schedulers; R-MCTS achieves the lowest count at every configuration
($16.3 \to 11.0 \to 1.4 \to 0.0$ as $N_c$ increases).
Table~\ref{tab:multicore} gives complete numerical results.

\begin{figure*}[!t]
  \centering
  \includegraphics[width=\linewidth]{fig12_multicore_scalability_2col.png}
  \caption{Multi-core scalability (Exp.\ 8): all schedulers vs core count
    $N_c{=}1,2,4,8$ on mixed workload.
    (a)~turnaround; (b)~starvation count; (c)~CPU utilisation.
    R-MCTS trained on $N_c{=}1$ only.}
  \label{fig:multicore}
\end{figure*}

\begin{table*}[!t]
\centering
\caption{Multi-Core Scalability: Mean Turnaround ($\pm$ SD), Starvation,
  CPU Utilisation, and Load Imbalance ($\sigma_U$) per Core Count
  (15 Trials, 20 Processes, Mixed Workload)}
\label{tab:multicore}
\renewcommand{\arraystretch}{1.1}
\begin{tabular}{@{}l r r r r r r r r@{}}
\toprule
& \multicolumn{2}{c}{$N_c{=}1$}
& \multicolumn{2}{c}{$N_c{=}2$}
& \multicolumn{2}{c}{$N_c{=}4$}
& \multicolumn{2}{c}{$N_c{=}8$} \\
\cmidrule(lr){2-3}\cmidrule(lr){4-5}\cmidrule(lr){6-7}\cmidrule(lr){8-9}
\textbf{Scheduler} & Turn. & Starv. & Turn. & $\sigma_U$ & Turn. & $\sigma_U$ & Turn. & $\sigma_U$ \\
\midrule
FCFS        & 167.59 & 16.2 & 82.02 & 0.031 & 41.54 & 0.074 & 26.92 & 0.153 \\
SJF-True    & 135.85 & 14.9 & 69.28 & 0.027 & 38.16 & 0.081 & 26.92 & 0.160 \\
SJF-Est     & 218.77 & 18.1 & 102.36& 0.019 & 46.33 & 0.051 & 27.08 & 0.161 \\
Round-Robin & 176.30 & 16.7 & 78.65 & 0.039 & 37.24 & 0.092 & 22.27 & 0.196 \\
\textbf{R-MCTS} & 182.39 & \textbf{16.3} & 88.56 & \textbf{0.024} & 41.62 & \textbf{0.062} & 27.01 & \textbf{0.157} \\
\bottomrule
\end{tabular}
\end{table*}

\subsection{Core-Count Transfer (Exp.\ 9)}
\label{sec:exp9}

Figure~\ref{fig:generalisation} shows R-MCTS performance relative to
FCFS on $N_c{=}1,2,4,8$ when the agent is trained \emph{exclusively}
on $N_c{=}1$ trajectories.
R-MCTS achieves statistically indistinguishable turnaround from FCFS at
$N_c{=}2$ ($p{=}0.995$), $N_c{=}4$ ($p{=}0.701$), and
$N_c{=}8$ ($p{=}0.943$), confirming zero-shot core-count transfer.
At $N_c{=}4$ R-MCTS shows a small positive improvement ($+2.2\%$),
suggesting the priority-aware value estimator provides a mild benefit
when four cores compete for high-priority processes.

The theoretical basis for this transfer property is that
Algorithm~\ref{alg:policy} is a per-process function: its inputs are
$f(p)$ (6 per-process features) and $\mathbf{s}$ (8 queue-level
features), neither of which encodes $N_c$.
The same policy is therefore valid for any core count at deployment
time, enabling zero-shot transfer without architecture changes.

\begin{figure}[!t]
  \centering
  \includegraphics[width=\linewidth]{fig13_generalisation_1col.png}
  \caption{Core-count transfer (Exp.\ 9):
    (a)~R-MCTS vs FCFS turnaround as core count increases;
    (b)~R-MCTS improvement over FCFS at each $N_c$.
    R-MCTS trained on $N_c{=}1$ only. All differences non-significant.}
  \label{fig:generalisation}
\end{figure}

\subsection{Google Trace Multi-Core Validation (Exp.\ 10)}

Figure~\ref{fig:googlemulti} shows R-MCTS on the Google-trace-derived
workload at $N_c{=}1,2,4$.
At $N_c{=}2$, R-MCTS achieves the lowest load imbalance
($\sigma_U{=}0.025$ vs FCFS 0.045, Round-Robin 0.060) and the lowest
starvation count (10.2 vs FCFS 12.3).
At $N_c{=}4$, R-MCTS achieves the best realistic-baseline turnaround
(141.01 vs FCFS 151.04, SJF-Est 151.06) and lowest starvation (3.1
vs FCFS 6.1, Round-Robin 6.7).
These results confirm that the load-imbalance advantage observed on
synthetic workloads (Experiment~8) extends to real-world workload
statistics.

\begin{figure*}[!t]
  \centering
  \includegraphics[width=\linewidth]{fig14_google_multicore_2col.png}
  \caption{Google trace multi-core validation (Exp.\ 10, 15 trials):
    (a)~turnaround; (b)~CPU utilisation; (c)~starvation count.
    R-MCTS achieves lowest starvation and best load balance
    at $N_c{=}2$ and $N_c{=}4$.}
  \label{fig:googlemulti}
\end{figure*}

% ─────────────────────────────────────────────────────────────────────────────
\section{Discussion}
\label{sec:discussion}
% ─────────────────────────────────────────────────────────────────────────────

\subsection{Why R-MCTS Trails FCFS on Raw Turnaround}

R-MCTS pays one context-switch tick per preemption and occasionally
delays a shorter process to serve a higher-priority one.
The cost is a 7.7\%--23.8\% higher turnaround relative to FCFS
depending on workload type; the benefit is a 90.0\% reduction in
short-process wait time and the lowest starvation count across all
tested conditions.
Raw turnaround does not capture scheduling quality in a multi-objective
setting: a scheduler that minimises turnaround by ignoring priority and
permitting indefinite starvation would be entirely unsuitable for
production use.

\subsection{Load Imbalance as a Novel SMP Metric}

The load imbalance metric $\sigma_U$ is not reported in prior OS
scheduling papers. Our results show a consistent ordering:
Round-Robin generates the highest imbalance (its fixed-quantum
preemption creates bursty core release patterns),
while SJF-Est generates the lowest imbalance at $N_c{=}2$ because its
burst estimates produce more uniform per-process CPU times.
R-MCTS achieves lower imbalance than FCFS and Round-Robin at every
multi-core configuration, suggesting its priority-weighted selection
naturally distributes work more evenly across cores.

\subsection{Core-Count Transfer: Why It Works}

The core-count transfer property (Experiment~9) arises because the
value estimator encodes per-process quality, not global queue geometry.
When core~$i$ queries $\Pi$ independently of core~$j$, it observes the
shared ready queue after the previous assignment and makes a locally
optimal decision.
The effective sample size seen at runtime therefore scales with $N_c$
(each of $N_c$ cores makes independent selections per scheduling
round), but the policy itself is unchanged.
This is a fundamentally different mechanism from multi-agent RL
approaches~\cite{mao2016}, which require coordinated training across
agents. R-MCTS achieves multi-core coordination as an emergent
property of per-process feature independence.

\subsection{Threats to Validity}

Three threats apply.
First, the discrete-event simulator omits real kernel latencies, cache
effects, and NUMA penalties.
Second, the Google trace workload is generated from published statistics
rather than raw trace data, which may omit temporal correlations.
Third, the starvation threshold $\delta{=}25$ requires recalibration
at different workload scales.
Future work should evaluate R-MCTS on real production kernel traces
(Google cluster traces~\cite{reiss2011google}, Azure VM logs) and
implement the policy as a Linux scheduling class module.

% ─────────────────────────────────────────────────────────────────────────────
\section{Conclusion}
\label{sec:conclusion}
% ─────────────────────────────────────────────────────────────────────────────

This paper presented R-MCTS, an offline MCTS-based process scheduler
that trains from stored single-core failed trajectories and deploys
without modification on SMP systems with up to eight cores.
Ten experiments across synthetic workloads, a Google cluster trace
workload, and four SMP configurations establish three main findings.

First, on single-core mixed workloads R-MCTS achieves turnaround
statistically indistinguishable from the SJF-True theoretical ceiling
($p{=}0.667$, gap 1.9\%), outperforms all practical baselines by
23.8\%--42.3\% ($p{<}0.001$), and reduces starvation-related wait
time by 90.0\%.

Second, on multi-core systems R-MCTS achieves the lowest load imbalance
($\sigma_U$) at $N_c{=}2$ on both synthetic and Google trace workloads,
confirming that its priority-aware value estimator distributes work
more evenly across cores than arrival-order or preemptive baselines.

Third, R-MCTS trained exclusively on $N_c{=}1$ trajectories achieves
statistically indistinguishable performance from FCFS at $N_c{=}2,4,8$
without retraining ($p{>}0.70$ at all multi-core configurations),
demonstrating zero-shot core-count transfer as an emergent property of
its per-process feature representation.

Future directions include replacing the cosine-similarity lookup with a
learned state embedding, extending to NUMA-aware scheduling with
per-socket ready queues, and validating on real Linux scheduler traces
using the \texttt{sched\_class} interface.

\section*{CRediT Author Statement}

\textbf{Sanusi Mustapha Babansoro}: Conceptualization, Methodology,
Software, Formal analysis, Investigation, Data curation, Writing --
original draft, Visualization.
\textbf{Abrha Gebriye Embafresu}: Validation, Writing -- review \&
editing.

\section*{Funding}

This research did not receive any specific grant from funding agencies
in the public, commercial, or not-for-profit sectors.

\section*{Data Availability}

The simulation code, experiment scripts, and generated data supporting
the findings of this study are openly available on GitHub at
\url{https://github.com/Digitalmustiii/rmcts-scheduler}.
A preprint of this manuscript is available at
\url{https://doi.org/10.5281/zenodo.21317887}.

\section*{Acknowledgment}

The authors thank the faculty of the School of Computer Science and
Engineering at UESTC for foundational instruction in operating systems
that motivated this work.
During manuscript preparation, the authors used Claude (Anthropic)
for LaTeX formatting assistance; all content was reviewed, edited,
and verified by the authors, who take full responsibility for the
published article.

\bibliographystyle{IEEEtran}
\begin{thebibliography}{30}

\bibitem{rmcts_github}
S.\,M.~Babansoro and A.\,G.~Embafresu,
``R-MCTS: Zero-Shot Multi-Core Process Scheduling via Offline Monte Carlo
Tree Search with Contrastive Value Learning --- Source Code,''
GitHub, 2026. [Online]. Available:
\url{https://github.com/Digitalmustiii/rmcts-scheduler}

\bibitem{stallings}
W.~Stallings, \emph{Operating Systems: Internals and Design Principles},
9th~ed.\ Pearson, 2018.

\bibitem{silberschatz}
A.~Silberschatz, P.\,B.~Galvin, and G.~Gagne,
\emph{Operating System Concepts}, 10th~ed.\ Wiley, 2018.

\bibitem{tanenbaum}
A.\,S.~Tanenbaum and H.~Bos,
\emph{Modern Operating Systems}, 4th~ed.\ Prentice Hall, 2014.

\bibitem{sundqn}
X.~Sun \emph{et~al.}, ``Dynamic Operating System Scheduling Using Double
DQN: A Reinforcement Learning Approach to Task Optimization,'' in
\emph{Proc.\ 8th Int.\ Conf.\ Advanced Algorithms and Control Engineering
(ICAACE)}, IEEE, 2025.

\bibitem{hanlearning}
H.~Han \emph{et~al.}, ``Performance Improvement of Linux CPU Scheduler
Using Policy Gradient Reinforcement Learning,'' \emph{IEEE Access}, 2020.

\bibitem{linnos}
S.~Hao \emph{et~al.}, ``LinnOS: Predictability on Unpredictable Flash
Storage with a Light Neural Network,'' in \emph{Proc.\ USENIX OSDI}, 2020.

\bibitem{mao2016}
H.~Mao \emph{et~al.}, ``Resource Management with Deep Reinforcement
Learning,'' in \emph{Proc.\ 15th ACM Workshop on Hot Topics in Networks
(HotNets)}, pp.~50--56, 2016.

\bibitem{decima}
H.~Mao \emph{et~al.}, ``Learning Scheduling Algorithms for Data
Processing Clusters,'' in \emph{Proc.\ ACM SIGCOMM}, 2019.

\bibitem{rjoubccpe}
G.~Rjoub, J.~Bentahar, O.\,A.~Wahab, and A.\,S.~Bataineh,
``Deep and reinforcement learning for automated task scheduling
in large-scale cloud computing systems,''
\emph{Concurrency Comput.:\ Pract.\ Exp.}, vol.~33, no.~23, e5919, 2021.

\bibitem{barrettccpe}
E.~Barrett, E.~Howley, and J.~Duggan,
``Applying reinforcement learning towards automating resource allocation
and application scalability in the cloud,''
\emph{Concurrency Comput.:\ Pract.\ Exp.}, vol.~25, no.~12,
pp.~1656--1674, 2013.

\bibitem{pangccpe}
X.~Pang \emph{et~al.},
``Deep Reinforcement Learning for Cloud Manufacturing Task Scheduling
via Sparse-Reward Optimization,''
\emph{Concurrency Comput.:\ Pract.\ Exp.}, 2026, doi:10.1002/cpe.70680.

\bibitem{jalali2024}
Z.~Jalali Khalil Abadi, N.~Mansouri, and M.\,M.~Javidi,
``Deep reinforcement learning-based scheduling in distributed systems:
a critical review,''
\emph{Knowl.\ Inf.\ Syst.}, vol.~66, no.~10, pp.~5709--5782, 2024.

\bibitem{levine2020}
S.~Levine, A.~Kumar, G.~Tucker, and J.~Fu,
``Offline Reinforcement Learning: Tutorial, Review, and Perspectives
on Open Problems,'' arXiv:2005.01643, 2020.

\bibitem{qtransformer}
Y.~Chebotar \emph{et~al.},
``Q-Transformer: Scalable Offline Reinforcement Learning via
Autoregressive Q-Functions,''
in \emph{Proc.\ CoRL}, PMLR, 2023.

\bibitem{dpo}
R.~Rafailov \emph{et~al.},
``Direct Preference Optimization: Your Language Model is Secretly a
Reward Model,''
\emph{Adv.\ Neural Inf.\ Process.\ Syst.}, vol.~36,
pp.~53728--53741, 2023.

\bibitem{alphallm}
X.~Wang \emph{et~al.},
``Towards Self-Improvement of LLMs via MCTS: Leveraging Stepwise
Knowledge with Curriculum Preference Learning,''
arXiv:2410.06508, 2024.

\bibitem{uct}
L.~Kocsis and C.~Szepesv\'{a}ri,
``Bandit Based Monte-Carlo Planning,''
in \emph{Proc.\ ECML}, Springer, 2006.

\bibitem{mcts_survey}
C.\,B.~Browne \emph{et~al.},
``A Survey of Monte Carlo Tree Search Methods,''
\emph{IEEE Trans.\ Comput.\ Intell.\ AI Games},
vol.~4, no.~1, pp.~1--43, 2012.

\bibitem{alphago}
D.~Silver \emph{et~al.},
``Mastering the game of Go with deep neural networks and tree search,''
\emph{Nature}, vol.~529, no.~7587, pp.~484--489, 2016.

\bibitem{muzero}
J.~Schrittwieser \emph{et~al.},
``Mastering Atari, Go, chess and shogi by planning with a learned model,''
\emph{Nature}, vol.~588, no.~7839, pp.~604--609, 2020.

\bibitem{mctsep}
H.~Xu \emph{et~al.},
``MCTS-EP: Empowering Embodied Planning with Online Preference
Optimization,'' arXiv:2509.17116, 2025.

\bibitem{reiss2011google}
C.~Reiss, J.~Wilkes, and J.\,L.~Hellerstein,
``Google cluster-usage traces: format + schema,''
Google Inc., Tech.\ Rep., 2011.

\bibitem{minet2018}
P.~Minet, L.~Nataf, and I.~Flauzac,
``Analyzing traces from a Google data center,''
in \emph{Proc.\ IEEE ISCC}, 2018.

\bibitem{fusfs}
Y.~Fu \emph{et~al.},
``SFS: Smart OS Scheduling for Serverless Functions,''
in \emph{Proc.\ SC22}, IEEE, 2022.

\bibitem{gaodeep}
W.~Gao \emph{et~al.},
``Deep Learning Workload Scheduling in GPU Datacenters: Taxonomy,
Challenges and Vision,'' arXiv:2205.11913, 2022.

\bibitem{suttonbarto}
R.\,S.~Sutton and A.\,G.~Barto,
\emph{Reinforcement Learning: An Introduction}, 2nd~ed.\ MIT Press, 2018.

\bibitem{huangjsa}
J.~Huang \emph{et~al.},
``Energy optimization for deadline-constrained parallel applications
on multi-ECU embedded systems,''
\emph{J.\ Syst.\ Archit.}, vol.~132, 102739, 2022.

\bibitem{sobhanijsa}
H.~Sobhani \emph{et~al.},
``REALISM: Reliability-aware energy management in multi-level
mixed-criticality systems with service level degradation,''
\emph{J.\ Syst.\ Archit.}, vol.~117, 102090, 2021.

\end{thebibliography}

\end{document}
